\section{The CAP-G Method}
\label{sec:method}

\subsection{Asset Context Score}

Each host carries three context attributes assigned from fixed, publicly-grounded
category mixes. \emph{Asset criticality tier} $\in$ \{MISSION\_CRITICAL 1.00, HIGH 0.70,
MEDIUM 0.40, LOW 0.15\} with fleet shares $5/20/50/25\%$ (FIPS~199 impact levels);
\emph{network exposure zone} $\in$ \{INTERNET\_FACING 1.00, DMZ 0.75, INTERNAL 0.40,
ISOLATED 0.10\} with shares $8/12/65/15\%$; and \emph{data sensitivity} $\in$
\{CJIS 1.00, PII\_CUI 0.60, PUBLIC 0.20\} with shares $15/45/40\%$. The Asset Context
Score is the weighted blend
%
\begin{equation}
\mathrm{ACS}(h) = c_1\,\mathrm{crit}(h) + c_2\,\mathrm{zone}(h) + c_3\,\mathrm{sens}(h),
\end{equation}
%
with $c_1{=}0.5,\,c_2{=}0.3,\,c_3{=}0.2$ fixed from priors (criticality-dominant,
mirroring FIPS~199 system categorization) and \emph{not} tuned on any seed. A high-water-
mark variant $\mathrm{ACS_{hwm}}(h)=\max(\mathrm{crit},\mathrm{zone},\mathrm{sens})$,
matching FIPS~199's literal max rule, is reported as a robustness check.

The design rationale for the weighted blend is to keep the score interpretable and to make
criticality dominant by construction, consistent with how FIPS~199 categorization is used in
practice, where the impact level of the information a system handles is the primary driver of
its protection requirements. We chose to weight the three dimensions rather than take their
literal maximum because a weighted blend degrades gracefully: a host that is highly critical but
internally zoned and public-data-only still earns a substantial score from the criticality term
alone, whereas a strict maximum would treat a single internet-facing attribute as equivalent to
top criticality. The high-water-mark variant is retained as a robustness check precisely because
it is the literal reading of the FIPS~199 max rule, and reporting both lets a reader see that the
finding does not hinge on the particular aggregation. Fixing the weights from priors rather than
fitting them on data is a deliberate conservatism: any in-sample tuning of the context weights
would risk borrowing strength from the evaluation target, and the point of the study is to
measure what a defensible, untuned context layer buys, not the ceiling a fitted one could reach.

\subsection{Hygiene-augmented base scorer}

The context-blind base scorer assigns each pair the score
%
\begin{equation}
\begin{split}
S_{\mathrm{base}}(h,c) = {}&\alpha\,\mathrm{EPSS}(c) + \beta\,\mathrm{HRS}(h) \\
&+ \gamma\,\mathrm{KEV_{rec}}(c) + \delta\,\mathrm{EPSS}(c)\,\mathrm{HRS}(h),
\end{split}
\end{equation}
%
where $\mathrm{EPSS}(c)$ is the exploit probability~\cite{jacobs2021epss},
$\mathrm{KEV_{rec}}(c)$ an exponential decay on days since KEV entry, and $\mathrm{HRS}(h)$
a host hygiene risk score in $[0,1]$ formed from patch posture, directory exposure, and
telemetry freshness. The weights $(\alpha,\beta,\gamma,\delta){=}(0.7,0.5,0.1,0.2)$ and the
hygiene-dimension weights $(0.5,0.3,0.2)$ are fixed constants, not tuned here. This is the
hygiene-augmented scorer of prior work~\cite{paper4hygieneprio,paper1vulnprio}, reproduced here
unchanged so that CAP-G is evaluated as a clean increment on top of an established baseline
rather than against a weakened straw man; carrying the base scorer's own constants verbatim
ensures that any measured difference is attributable to asset context and not to a re-tuned
hygiene term. The multiplicative interaction term $\delta\,\mathrm{EPSS}(c)\,\mathrm{HRS}(h)$
is what lets a host that is both exposed-by-exploit and poorly maintained rise above either
signal alone, which is the behaviour a hygiene-augmented scorer is meant to capture. This is the
baseline against which CAP-G is measured.

\subsection{The CAP-G scorer}

CAP-G multiplies the base score by a host context factor:
%
\begin{equation}
S_{\mathrm{CAP\text{-}G}}(h,c) = S_{\mathrm{base}}(h,c)\cdot\bigl(1+\rho\,\mathrm{ACS}(h)\bigr),
\label{eq:capg}
\end{equation}
%
where $\rho\geq 0$ is the context-emphasis hyperparameter; $\rho{=}0$ recovers the base
scorer. The multiplicative form is a deliberate design choice over an additive one. Multiplying
keeps context as a modulator of the exploit-and-hygiene signal rather than a competing additive
term: a host with high mission value but no exploitable, poorly-maintained vulnerability still
contributes nothing to the queue, because its base score is near zero and the context factor
scales zero to zero. An additive context term, by contrast, could float a quiescent
high-criticality host to the top on context alone, which is not the behaviour a remediation queue
should have. The factor $1+\rho\,\mathrm{ACS}(h)$ also has a clean degenerate case: because it is
monotone in ACS, on a \emph{homogeneous} fleet (constant ACS) it is a constant rescaling that
leaves the ranking, and therefore every ranking metric, identical to the base scorer. That is not
an accident of the constants but a property of the form, and it yields an exact mechanism control
(Section~\ref{sec:results}) in which any non-zero advantage must come from fleet heterogeneity
rather than from the transform itself.

\subsection{Calibration}

$\rho$ was grid-searched on 5 held-out calibration seeds (objective: mean MWP@50), grid
$\{0.5,1,2,4,8\}$, and fixed before the 25-seed evaluation. MWP@50 rose monotonically
across the grid ($0.272\!\rightarrow\!0.356$), so the selected $\rho{=}8.0$ sits at the
grid boundary, with a $10.4$~pp calibration-set margin over the base scorer (well above the
$3$~pp stop-rule floor). The boundary saturation is reported transparently: a larger
$\rho$ would likely raise calibration MWP@50 further before degenerating toward the weak
context-only behaviour, implying an interior optimum beyond the locked grid. Per the
pre-registered stop rule we did \emph{not} expand the grid post hoc; the locked-grid
winner is used as-is, so the evaluation uses a conservative lower bound on context
emphasis.
